\SGAchapter{III}{相对有限性条件}{Conditions de finitude relatives}
\begin{center}L. Illusie\end{center}
\SourcePage{232}
\SGAsection{1}{相对伪凝聚性}
\begin{SGAresult}{\Term{theorem}{定理}}{1.1}
设 $n\in\mathbf Z$。对每个\Term{locally-finite-type}{局部有限型}的\Term{scheme}{概形}\Term{morphism}{态射} $f:X\to Y$，存在唯一的 $D(X)$\footnote{$D(X)$ 是 $\mathcal O_X$-\Term{module-category}{模范畴}的\Term{derived-category}{导出范畴}。}的\Term{full-subcategory}{满子范畴} $D(f)_{n\text{-coh}}$，满足下列条件：
\begin{enumerate}
\item[(i)] 若 $E\in\operatorname{Ob}D(f)_{n\text{-coh}}$，则对每个\Term{open-immersion}{开浸入} $i:U\to X$，有 $i^*E\in\operatorname{Ob}D(fi)_{n\text{-coh}}$。
\item[(ii)] 若 $E\in\operatorname{Ob}D(X)$，并且存在\Term{open-covering}{开覆盖} $(i_\alpha:U_\alpha\to X)_{\alpha\in A}$，使所有 $\alpha\in A$ 都满足 $i_\alpha^*E\in\operatorname{Ob}D(fi_\alpha)_{n\text{-coh}}$，则 $E\in\operatorname{Ob}D(f)_{n\text{-coh}}$。
\item[(iii)] 若
\[
\begin{tikzcd}X\arrow[r,"f"]\arrow[dr,"h"']&Y\arrow[d,"g"]\\&Z\end{tikzcd}
\]
是\Term{locally-finite-type}{局部有限型}\Term{morphism}{态射}的\Term{commutative-diagram}{交换图}，其中 $f$ 是\Term{closed-immersion}{闭浸入}，$g$ 是\Term{smooth}{光滑态射}，则
\[
E\in\operatorname{Ob}D(h)_{n\text{-coh}}
\quad\Longleftrightarrow\quad f_*E\in D(Y)_{n\text{-coh}}.
\]
\footnote{$D(Y)_{n\text{-coh}}$ 是由 $n$-\Term{pseudo-coherent}{伪凝聚}\Term{object}{对象}组成的 $D(Y)$ 的\Term{full-subcategory}{满子范畴}（第 I 讲 2.15）。}
\end{enumerate}
\end{SGAresult}
\begin{proof}
唯一性显然。事实上，设 $f:X\to Y$ 为\Term{locally-finite-type}{局部有限型}\Term{morphism}{态射}。取\Term{open-covering}{开覆盖} $(i_\alpha:X_\alpha\to X)_{\alpha\in A}$，使每个 $fi_\alpha$ 均\Term{decomposable}{可分解}为
\[
X_\alpha\xrightarrow{i'_\alpha}Y_\alpha\xrightarrow{f'_\alpha}Y,
\]
其中 $i'_\alpha$ 是\Term{closed-immersion}{闭浸入}，$f'_\alpha$ 是\Term{smooth}{光滑态射}。由 (i)、(ii)、(iii) 立即得知，$D(f)_{n\text{-coh}}$ 的\Term{object}{对象}正是那些满足 $i'_{\alpha*}i_\alpha^*E\in\operatorname{Ob}D(Y_\alpha)_{n\text{-coh}}$ 的 $E\in\operatorname{Ob}D(X)$。
\SourcePage{233}
因此 $D(f)_{n\text{-coh}}$ 确实被唯一确定，因为它是 $D(X)$ 的\Term{full-subcategory}{满子范畴}。为证明存在性，我们需要几个\Term{lemma}{引理}。
\end{proof}
\begin{SGAresult}{\Term{lemma}{引理}}{1.1.1}
设 $i:X\to Y$ 是\Term{scheme}{概形}的\Term{closed-immersion}{闭浸入}，且 $i_*\mathcal O_X$ \Term{pseudo-coherent}{伪凝聚}；设 $E\in\operatorname{Ob}D(X)$。则 $E$ 是 $n$-\Term{pseudo-coherent}{伪凝聚}的，当且仅当 $i_*E$ 是如此。
\end{SGAresult}
\begin{proof}
必要性来自\Term{pseudo-coherent}{伪凝聚性}在施行\Term{functor}{函子}后保持\Term{stable}{稳定}的\Term{generic}{一般}判据（第 I 讲 2.13）。反过来，设 $i_*E$ 是 $n$-\Term{pseudo-coherent}{伪凝聚}的，证明 $E$ 也如此。由于 $i_*E$ 的\Term{cohomology}{上同调}\Term{local-scheme}{局部}有上界，$E$ 也如此；问题是\Term{local-scheme}{局部}的，因此可假定对充分大的 $m$，$E$ 是 $m$-\Term{pseudo-coherent}{伪凝聚}的。对 $m$ 作\Term{descent}{下降}归纳，证明对每个 $m\ge n$，$E$ 都是 $m$-\Term{pseudo-coherent}{伪凝聚}的。可假定已知 $E$ 是 $(n+1)$-\Term{pseudo-coherent}{伪凝聚}的。\Term{localization}{局部化}后，有\Term{distinguished-triangle}{可区别三角形}
\[
E'\longrightarrow E\longrightarrow E''\longrightarrow E'[1],
\]
其中 $E'$ \Term{perfect-complex}{完美}，且 $E''{}^i=0$（$i>n$）。由第 I 讲 2.6，只须证明 $E''$ 是 $n$-\Term{pseudo-coherent}{伪凝聚}的。对上述三角形施行 $i_*$，得到\Term{distinguished-triangle}{可区别三角形}
\[
i_*E'\longrightarrow i_*E\longrightarrow i_*E''\longrightarrow(i_*E')[1].
\]
由证明的前半部分，$i_*E'$ \Term{pseudo-coherent}{伪凝聚}；故由第 I 讲 2.6，$i_*E''$ 是 $n$-\Term{pseudo-coherent}{伪凝聚}的。由于 $i_*E''$ 在\Term{degree}{次数} $>n$ 为零，由第 I 讲 2.10 b)，这意味着 $H^n(i_*E'')$ 为\Term{finite-type}{有限型}。又因为 $H^n(i_*E'')\simeq i_*H^n(E'')$，再次由第 I 讲 2.10 b)，归结为证明：若一个 $\mathcal O_X$-模 $M$ 的 $i_*M$ 是\Term{finite-type}{有限型}的，则 $M$ 也是\Term{finite-type}{有限型}的。这是直接的，而且不依赖于关于 $i_*\mathcal O_X$ 的假设。
\end{proof}
\SourcePage{234}
\textbf{注记 1.1.2。} 上一\Term{lemma}{引理}特别适用于 $i$ 为\Term{regular-immersion}{正则浸入}的情形（第 VII 讲 1.4）。事实上，借助 Koszul \Term{complex}{复形}，\Term{local-scheme}{局部}存在 $i_*\mathcal O_X$ 的一个由\Term{finite-type}{有限型}自由 $\mathcal O_Y$-模组成的\Term{finite}{有限}左\Term{resolution}{解消}，因此 $i_*\mathcal O_X$ \Term{perfect-complex}{完美}（第 I 讲 4.8、4.14），从而尤其\Term{pseudo-coherent}{伪凝聚}。
\begin{SGAresult}{\Term{lemma}{引理}}{1.1.3}
设有\Term{scheme}{概形}的\Term{cartesian-square}{笛卡尔方块}
\[
\begin{tikzcd}X\arrow[d,"i"']&X'\arrow[l,"f"']\arrow[d,"i'"]\\Y&Y'\arrow[l,"g"]\end{tikzcd}
\]
其中 $i$ 是\Term{closed-immersion}{闭浸入}，并设 $E$ 为 $\mathcal O_X$-模\Term{complex}{复形}。则：
\begin{enumerate}
\item[(i)] \Term{canonical}{典范}的\Term{base-change}{基变换}\Term{arrow}{箭头}（EGA III 1.4.15.4）
\[g^*i_*E\longrightarrow i'_*f^*E\]
是\Term{isomorphism}{同构}。
\item[(ii)] 令 $E'=f^*E$。若 $g$ \Term{flat}{平坦}且 $i_*E$ 是 $n$-\Term{pseudo-coherent}{伪凝聚}的，则 $i'_*E'$ 是 $n$-\Term{pseudo-coherent}{伪凝聚}的。反过来，若 $g$ \Term{faithfully-flat}{忠实平坦}且 $i'_*E'$ 是 $n$-\Term{pseudo-coherent}{伪凝聚}的，则 $i_*E$ 是 $n$-\Term{pseudo-coherent}{伪凝聚}的。
\end{enumerate}
\end{SGAresult}
\begin{proof}该断言在\Term{stalk}{茎}上直接可验。\end{proof}
\SourcePage{235}
\begin{SGAresult}{\Term{lemma}{引理}}{1.1.4}
设有\Term{scheme}{概形}的\Term{commutative-diagram}{交换图}
\[
\begin{tikzcd}X'\arrow[dr,"f'"']&X\arrow[l,"i'"']\arrow[r,"i''"]&X''\arrow[dl,"f''"]\\&Y&\end{tikzcd}
\]
其中 $i'$（相应地 $i''$）为\Term{closed-immersion}{闭浸入}，$f'$（相应地 $f''$）为\Term{smooth}{光滑态射}，并设 $E\in\operatorname{Ob}D(X)$。则
\[
i'_*E\text{ 是 }n\text{-伪凝聚的}\quad\Longleftrightarrow\quad i''_*E\text{ 是 }n\text{-伪凝聚的}.
\]
\end{SGAresult}
\begin{proof}
采用 Lichtenbaum 的一个技巧，将 $X$ 对角地映入 $X'\times_YX''$，便立即归结为 $f'$、$f''$ 中一个为恒等\Term{morphism}{态射}的情形。改换记号，有\Term{commutative-diagram}{交换图}
\[
\begin{tikzcd}&Y\arrow[d,"f"]\\X\arrow[ur,"i"]\arrow[r,"j"']&Z\end{tikzcd}\tag{*}
\]
其中 $i,j$ 为\Term{closed-immersion}{闭浸入}，$f$ 为\Term{smooth}{光滑态射}。须证明 $i_*E$ 是 $n$-\Term{pseudo-coherent}{伪凝聚}的\Term{equivalence}{等价}于 $j_*E$ 是 $n$-\Term{pseudo-coherent}{伪凝聚}的。此问题在 $Y$ 上以及在 $Z$ 中 $X$ 的某点附近都是\Term{local-scheme}{局部}的，因此可假定 $X,Y,Z$ 均\Term{affine}{仿射}，且存在\Term{resolution}{分解}
\[
\begin{tikzcd}Y\arrow[r,"i'"]\arrow[d,"f"']&Z[T_1,\ldots,T_r]\arrow[dl,"f'"]\\Z&\end{tikzcd}
\]
其中 $i'$ 是\Term{closed-immersion}{闭浸入}，斜\Term{arrow}{箭头}是\Term{canonical-projection}{典范投影}。此时 $i'$ 是\Term{regular-immersion}{正则浸入}（第 VII 讲 1.10），故由 1.1.1，$D(Y)$ 的\Term{object}{对象} $F$ 是 $n$-\Term{pseudo-coherent}{伪凝聚}的，当且仅当 $i'_*F$ 如此。
\SourcePage{236}
于是可假定 $Y=Z[T_1,\ldots,T_r]$，$f$ 为\Term{canonical-projection}{典范投影}。图 (*) 对应于环图
\[
\begin{tikzcd}&B=C[T_1,\ldots,T_r]\arrow[dl,"u"']\\A&C\arrow[l,"v"]\arrow[u,"e"']\end{tikzcd}
\]
其中 $u,v$ 为\Term{epimorphism}{满态射}，$e$ 是\Term{canonical}{典范}包含。对 $1\le i\le r$，令 $u(T_i)=x_i$，取 $z_i\in C$ 满足 $v(z_i)=x_i$。由 $\sigma(T_i)=z_i$ 定义的 $e$ 的\Term{retraction}{收缩} $\sigma:B\to C$ 满足 $v\sigma=u$。用几何语言来说，它对应于 $f$ 的一个\Term{section}{截面} $s$，且 $i=sj$。由第 VII 讲 1.10，$s$ 是\Term{regular-immersion}{正则浸入}，所以断言来自 1.1.1。
\end{proof}
\begin{proof}[\Term{theorem}{定理} 1.1 的证明（续）]
设 $f:X\to Y$ 为\Term{locally-finite-type}{局部有限型}\Term{morphism}{态射}。存在 $X$ 的\Term{open-covering}{开覆盖} $(X_\alpha)$，使每个 $X_\alpha$ 以\Term{closed-immersion}{闭浸入}\Term{embedding}{嵌入}一个在 $Y$ 上\Term{smooth}{光滑}的\Term{scheme}{概形} $Y_\alpha$。令 $X'=\coprod X_\alpha$、$Y'=\coprod Y_\alpha$，便得到\Term{commutative-diagram}{交换图}
\[
\begin{tikzcd}X'\arrow[r,"i'"]\arrow[d,"p'"']&Y'\arrow[ddl,"f'"]\\X\arrow[d,"f"']&\\Y&\end{tikzcd}\tag{a'}
\]
其中 $p'$ 是一族\Term{covering}{覆盖}其目标的\Term{open-immersion}{开浸入}之和（为简便起见，称之为\Term{special-morphism}{特殊态射}），$i'$ 是\Term{closed-immersion}{闭浸入}，$f'$ 是\Term{smooth}{光滑态射}。定义 $D(f)_{n\text{-coh}}$ 为由满足 $i'_*p'^*E$ 是 $n$-\Term{pseudo-coherent}{伪凝聚}的\Term{object}{对象} $E$ 组成的 $D(X)$ 的\Term{full-subcategory}{满子范畴}。先证明此定义有意义，即它不依赖 $i',p',f'$ 的选择。
\SourcePage{237}
设另有类似于 (a') 的\Term{commutative-diagram}{交换图}
\[
\begin{tikzcd}X''\arrow[r,"i''"]\arrow[d,"p''"']&Y''\arrow[ddl,"f''"]\\X\arrow[d,"f"']&\\Y&\end{tikzcd}\tag{a''}
\]
其中 $p''$ 为\Term{special-morphism}{特殊态射}，$i''$ 为\Term{closed-immersion}{闭浸入}，$f''$ 为\Term{smooth}{光滑态射}。由 (a')、(a'') 得到\Term{commutative-diagram}{交换图}
\[
\begin{tikzcd}[column sep=large,row sep=large]
X'\times_XX''\arrow[rr]\arrow[dr]\arrow[dd]& &Y'\times_YY''\arrow[dr]\arrow[dd]&\\
&X'\arrow[rr,"i'"]\arrow[dd,"p'"']&&Y'\arrow[dd,"f'"]\\
X''\arrow[rr,"i''" near start]\arrow[dr,"p''"']&&Y''\arrow[dr,"f''"']&\\
&X\arrow[rr,"f"']&&Y
\end{tikzcd}
\]
其中除 $f$ 外的各水平\Term{arrow}{箭头}均为\Term{closed-immersion}{闭浸入}；左侧\Term{cartesian-square}{笛卡尔方块}中的\Term{arrow}{箭头}为\Term{special-morphism}{特殊态射}，右侧的为\Term{smooth}{光滑态射}。因此只须证明如下断言。

\textbf{(A)} 设
\[
\begin{tikzcd}X_1\arrow[r,"i_1"]\arrow[d,"p"']&Y_1\arrow[d,"q"]\\X_2\arrow[r,"i_2"']&Y_2\end{tikzcd}
\]
是\Term{scheme}{概形}的\Term{commutative-diagram}{交换方块}，其中 $p$ 为\Term{special-morphism}{特殊态射}，$q$ 为\Term{smooth}{光滑态射}，$i_1,i_2$ 为\Term{closed-immersion}{闭浸入}；设 $E\in\operatorname{Ob}D(X_2)$。则 $i_{2*}E$ 是 $n$-\Term{pseudo-coherent}{伪凝聚}的，当且仅当 $i_{1*}p^*E$ 如此。
\SourcePage{238}
\textbf{(A) 的证明。} 按定义，$p$ 是 $X_2$ 的一族\Term{open-subset}{开集} $U_\alpha$ 之和，这族\Term{open-subset}{开集}\Term{covering}{覆盖} $X_2$。每个 $U_\alpha$ 由 $Y_2$ 的某个\Term{open-subset}{开集} $V_\alpha$ 诱导，而 $V_\alpha$ \Term{covering}{覆盖} $Y_2$ 中 $X_2$ 的一个\Term{open-neighborhood}{开邻域} $X'_2$。于是有\Term{commutative-diagram}{交换图}，其中方块为\Term{cartesian-square}{笛卡尔方块}：
\[
\begin{tikzcd}X_1\arrow[r,"i'_1"]\arrow[d,"p"']&X'_1\arrow[d,"p'"]&\\X_2\arrow[r,"i'_2"']\arrow[rr,bend right=20,"i_2"']&X'_2\arrow[r,"j"']&Y_2\end{tikzcd}
\]
这里 $i'_1,i'_2$ 为\Term{closed-immersion}{闭浸入}，$p,p'$ 为\Term{special-morphism}{特殊态射}，$j$ 为\Term{open-immersion}{开浸入}。$i_{2*}E$ 是 $n$-\Term{pseudo-coherent}{伪凝聚}的，显然\Term{equivalence}{等价}于 $i'_{2*}E$ 是 $n$-\Term{pseudo-coherent}{伪凝聚}的。另一方面，由 1.1.3，后者\Term{equivalence}{等价}于 $i'_{1*}p^*E$ 是 $n$-\Term{pseudo-coherent}{伪凝聚}的。再对 $i_2p$ 的两个\Term{resolution}{分解} $(jp')i'_1$、$qi_1$ 以及 $p^*E\in\operatorname{Ob}D(X_1)$ 应用 1.1.4，最后得到 $i_{2*}E$ 是 $n$-\Term{pseudo-coherent}{伪凝聚}的当且仅当 $i_{1*}p^*E$ 如此，证毕。

因此，通过图 (a') 定义的\Term{category}{范畴} $D(f)_{n\text{-coh}}$ 实际上只依赖于 $f$。此外，从定义直接可见它满足 1.1 的条件 (i)--(iii)。\Term{theorem}{定理} 1.1 证毕。
\end{proof}
\begin{SGAresult}{定义}{1.2}
设 $f:X\to Y$ 为\Term{locally-finite-type}{局部有限型}\Term{morphism}{态射}，$E\in\operatorname{Ob}D(X)$，$n\in\mathbf Z$。若 $E\in\operatorname{Ob}D(f)_{n\text{-coh}}$（见 1.1），称 $E$ \Term{relative-pseudo-coherent}{相对于 $f$ 为 $n$-伪凝聚}。若对所有 $n\in\mathbf Z$ 都如此，称 $E$ \Term{relative-pseudo-coherent}{相对于 $f$ 伪凝聚}（也称 $(-\infty)$-\Term{relative-pseudo-coherent}{相对伪凝聚}）。记 $D(f)_{\mathrm{coh}}$ 为由\Term{relative-pseudo-coherent}{相对于 $f$ 伪凝聚}的\Term{object}{对象}组成的 $D(X)$ 的\Term{full-subcategory}{满子范畴}。
\SourcePage{239}
若 $\mathcal O_X$ \Term{relative-pseudo-coherent}{相对于 $f$ 为 $n$-伪凝聚}（相应地，\Term{relative-pseudo-coherent}{伪凝聚}），便称 $f$ 为 $n$-\Term{pseudo-coherent}{伪凝聚}（相应地，\Term{pseudo-coherent}{伪凝聚}）\Term{morphism}{态射}。不致引起混淆时，也说“相对于 $Y$”，并以 $D(X)_{Y\text{-}n\text{-coh}}$（相应地 $D(X)_{Y\text{-coh}}$）代替 $D(f)_{n\text{-coh}}$（相应地 $D(f)_{\mathrm{coh}}$）。最后，对 $*=+,-,b$，令
\[
D^*(f)_{n\text{-coh}}=D^*(X)\cap D(f)_{n\text{-coh}},\qquad
D^*(f)_{\mathrm{coh}}=D^*(X)\cap D(f)_{\mathrm{coh}}.
\]
\end{SGAresult}
\textbf{注记 1.2.1。} a) \Term{category}{范畴} $D(f)_{n\text{-coh}}$、$D(f)_{\mathrm{coh}}$ \Term{generic}{一般}并不包含于 $D^-(X)$。不过当 $X$ \Term{quasi-compact}{拟紧}时却是如此，因为此时一个\Term{cohomology}{上同调}\Term{local-scheme}{局部}有上界的 $D(X)$ \Term{object}{对象}，其\Term{cohomology}{上同调}便有上界。

b) $D(f)_{\mathrm{coh}}$ 的\Term{object}{对象}具有\Term{quasi-coherent}{拟凝聚}\Term{cohomology}{上同调}。事实上，设 $E\in\operatorname{Ob}D(f)_{\mathrm{coh}}$。问题是\Term{local-scheme}{局部}的，可假定 $f$ 是\Term{closed-immersion}{闭浸入}。此时 $f_*E$ \Term{pseudo-coherent}{伪凝聚}，因而具有\Term{quasi-coherent}{拟凝聚}\Term{cohomology}{上同调}；再利用一个 $\mathcal O_X$-模 $M$ \Term{quasi-coherent}{拟凝聚}当且仅当 $f_*M$ 如此即可。

c) 对 $\mathcal O_X$-模，我们采用“\Term{relative-finite-presentation}{相对于 $f$ 有限表示}”这一术语来代替“\Term{relative-pseudo-coherent}{相对于 $f$ 为 $(-1)$-伪凝聚}”。由第 I 讲 2.9，一个 $\mathcal O_X$-模 $E$ \Term{relative-finite-presentation}{相对于 $f$ 有限表示}，当且仅当\Term{local-scheme}{局部}存在\Term{closed-immersion}{闭浸入} $i:X\to X'$，其中 $X'$ 在 $Y$ 上\Term{smooth}{光滑}，使 $i_*E$（按通常意义）\Term{finite-presentation}{有限表示}。特别地，
\SourcePage{240}
$f$ 是 $(-1)$-\Term{pseudo-coherent}{伪凝聚}的\Term{equivalence}{等价}于 $f$ 按 EGA IV 1.4 的意义\Term{locally-finite-presentation}{局部有限表示}。

\textbf{例 1.3。} \Term{smooth}{光滑态射}和\Term{regular-immersion}{正则浸入}（1.1.2）都是\Term{pseudo-coherent}{伪凝聚}\Term{morphism}{态射}。若 $Y$ \Term{locally-noetherian}{局部 Noether}，任意\Term{locally-finite-type}{局部有限型}\Term{morphism}{态射} $f:X\to Y$ 也如此。
\begin{SGAresult}{命题}{1.4}
设 $f:X\to Y$ 是\Term{locally-finite-type}{局部有限型}\Term{morphism}{态射}，$E\to F\to G\to E[1]$ 为 $D(X)$ 的\Term{distinguished-triangle}{可区别三角形}。则下表成立；其中 $n\in\mathbf Z\cup\{-\infty\}$ 表示\Term{relative-pseudo-coherent}{相对于 $f$ 的伪凝聚}\Term{degree}{次数}，每一行表示一个蕴含关系，下划线标出其结论：
\[
\begin{array}{ccc}E&F&G\\n+1&n&\underline n\\n&\underline n&n\\\underline n&n&n-1\end{array}.
\]
特别地，$D(f)_{\mathrm{coh}}$ 是 $D(X)$ 的\Term{triangulated-subcategory}{三角子范畴}。
\end{SGAresult}
\begin{proof}
由 1.1 的性质 (i)--(iii)，可假定 $f$ 为\Term{closed-immersion}{闭浸入}。于是有\Term{distinguished-triangle}{可区别三角形}
\[f_*E\longrightarrow f_*F\longrightarrow f_*G\longrightarrow(f_*E)[1],\]
断言遂由第 I 讲 2.6 得到。
\end{proof}
\begin{SGAresult}{定义}{1.5}
设 $X,Y$ 为两个 $S$-\Term{scheme}{概形}。若
\[\operatorname{Tor}_i^S(\mathcal O_X,\mathcal O_Y)=0\qquad(i\ne0),\]
称 $X,Y$ 在 $S$ 上\Term{tor-independent}{Tor 无关}；这里 $\operatorname{Tor}_i^S$ 是 EGA III 6.5 定义的\Term{local-scheme}{局部}超 Tor。
\end{SGAresult}
\SourcePage{241}
这是一个\Term{local-scheme}{局部}条件；若 $X,Y,S$ 是分别对应于环 $B,C,A$ 的\Term{affine}{仿射}\Term{scheme}{概形}，此条件意味着\Term{canonical}{典范}\Term{arrow}{箭头}
\[B\mathbin{\overset{\mathrm L}{\otimes}}_AC\longrightarrow B\otimes_AC\]
为\Term{isomorphism}{同构}。因此，特别当 $X$ 或 $Y$ 在 $S$ 上\Term{flat}{平坦}时，它成立。
\begin{SGAresult}{\Term{lemma}{引理}}{1.5.1}
设有\Term{scheme}{概形}的\Term{cartesian-square}{笛卡尔方块}
\[
\begin{tikzcd}X\arrow[d]&X'\arrow[l]\arrow[d]&X''\arrow[l]\arrow[d]\\Y&Y'\arrow[l]&Y''\arrow[l]\end{tikzcd}
\]
若 $X,Y'$ 在 $Y$ 上\Term{tor-independent}{Tor 无关}，且 $X',Y''$ 在 $Y'$ 上\Term{tor-independent}{Tor 无关}，则 $X,Y''$ 在 $Y$ 上\Term{tor-independent}{Tor 无关}。
\end{SGAresult}
\begin{proof}这是关于\Term{tensor-product}{张量积}\Term{associative-law}{结合律}的一个简单练习。\end{proof}
\textbf{记号 1.6。} 设 $X,Y$ 为两个 $S$-\Term{scheme}{概形}，$E\in\operatorname{Ob}D^-(X)$，$F\in\operatorname{Ob}D^-(Y)$。令
\[E\mathbin{\overset{\mathrm L}{\otimes}}_SF
=\mathrm L\operatorname{pr}_1^*E\mathbin{\overset{\mathrm L}{\otimes}}\mathrm L\operatorname{pr}_2^*F,\]
其中 $\operatorname{pr}_1,\operatorname{pr}_2$ 为 $X\times_SY$ 的两个\Term{projection}{投影}。
\begin{SGAresult}{\Term{lemma}{引理}}{1.7}
设 $S$ 为\Term{scheme}{概形}；对 $i=1,2$，设 $f_i:X_i\to Y_i$ 为 $S$-\Term{closed-immersion}{闭浸入}。假定 $X_1,X_2$ 以及 $Y_1,Y_2$ 在 $S$ 上\Term{tor-independent}{Tor 无关}。则对 $E_i\in\operatorname{Ob}D^-(X_i)_{\mathrm{qcoh}}$\footnote{$D(X_i)_{\mathrm{qcoh}}$ 表示由具有\Term{quasi-coherent}{拟凝聚}\Term{cohomology}{上同调}的\Term{complex}{复形}组成的 $D(X_i)$ 的\Term{full-subcategory}{满子范畴}。}，存在\Term{canonical}{典范}\Term{isomorphism}{同构}
\[
f_{1*}E_1\mathbin{\overset{\mathrm L}{\otimes}}_Sf_{2*}E_2
\xrightarrow{\sim}(f_1\times_Sf_2)_*(E_1\mathbin{\overset{\mathrm L}{\otimes}}_SE_2).
\]
\end{SGAresult}
\begin{proof}
令 $f=f_1\times_Sf_2$。由\Term{adjunction}{伴随}容易定义一个\Term{canonical}{典范}\Term{arrow}{箭头}
\SourcePage{242}
\[f_{1*}E_1\mathbin{\overset{\mathrm L}{\otimes}}_Sf_{2*}E_2\longrightarrow f_*(E_1\mathbin{\overset{\mathrm L}{\otimes}}_SE_2).\]
证明它为\Term{isomorphism}{同构}。问题是\Term{local-scheme}{局部}的，可假定 $S,X_i,Y_i$ 是分别对应于环 $C,A_i,B_i$ 的\Term{affine}{仿射}\Term{scheme}{概形}。逐步归约，可归结为 $E_i=\mathcal O_{X_i}$。此时须验证
\[A_1\mathbin{\overset{\mathrm L}{\otimes}}_{B_1}B\mathbin{\overset{\mathrm L}{\otimes}}_{B_2}A_2\longrightarrow A\]
是\Term{isomorphism}{同构}，其中 $A=A_1\otimes_CA_2$，$B=B_1\otimes_CB_2$。结合\Term{tor-independent}{Tor 无关}的假设，这由\Term{tensor-product}{张量积}的\Term{associative-law}{结合律}立即得到。
\end{proof}
\begin{SGAresult}{命题}{1.8}
设 $S$ 为\Term{scheme}{概形}；对 $i=1,2$，设 $f_i:X_i\to Y_i$ 为\Term{locally-finite-type}{局部有限型}的 $S$-\Term{morphism}{态射}。假定 $X_1,X_2$ 以及 $Y_1,Y_2$ 在 $S$ 上\Term{tor-independent}{Tor 无关}。设 $[a_i,b_i]$ 为 $\mathbf Z$ 中的区间，$E_i\in\operatorname{Ob}D^-(X_i)_{\mathrm{qcoh}}$ 是\Term{relative-pseudo-coherent}{相对于 $f_i$ 为 $a_i$-伪凝聚}且在\Term{degree}{次数} $>b_i$ \Term{acyclic}{无上同调}的\Term{complex}{复形}。则 $E_1\mathbin{\overset{\mathrm L}{\otimes}}_SE_2$ \Term{relative-pseudo-coherent}{相对于 $f=f_1\times_Sf_2$ 为 $a$-伪凝聚}，且在\Term{degree}{次数} $>b$ \Term{acyclic}{无上同调}，其中
\[a=\sup(a_1+b_2,a_2+b_1),\qquad b=b_1+b_2.\]
\end{SGAresult}
\begin{proof}
问题在 $X_1\times_SX_2$ 上是\Term{local-scheme}{局部}的，可假定 $f_i$ \Term{resolution}{分解}为 $X_i\xrightarrow{f'_i}Y'_i\xrightarrow{p_i}Y_i$，其中 $f'_i$ 为\Term{closed-immersion}{闭浸入}，$p_i$ 为\Term{smooth}{光滑态射}。令 $f'=f'_1\times_Sf'_2$，须证明 $f'_*(E_1\mathbin{\overset{\mathrm L}{\otimes}}_SE_2)$ 为 $a$-\Term{pseudo-coherent}{伪凝聚}且在\Term{degree}{次数} $>b$ \Term{acyclic}{无上同调}。由 1.5.1，$Y'_1,Y'_2$ 在 $S$ 上\Term{tor-independent}{Tor 无关}；故改换记号后，可假定 $f_i=f'_i$，即 $f_i$ 为\Term{closed-immersion}{闭浸入}。按假设，$f_{i*}E_i$ 为 $a_i$-\Term{pseudo-coherent}{伪凝聚}且在\Term{degree}{次数} $>b_i$ \Term{acyclic}{无上同调}；故由第 I 讲 2.16，$f_{1*}E_1\mathbin{\overset{\mathrm L}{\otimes}}_Sf_{2*}E_2$ 为 $a$-\Term{pseudo-coherent}{伪凝聚}且在\Term{degree}{次数} $>b$ \Term{acyclic}{无上同调}。再由 1.7 即得结论。
\end{proof}
\SourcePage{243}
\begin{SGAresult}{推论}{1.9}
在 1.8 的假设下，\Term{functor}{函子}
\[\mathbin{\overset{\mathrm L}{\otimes}}_S:D^-(X_1)\times D^-(X_2)\longrightarrow D^-(X)\]
诱导一个\Term{functor}{函子}
\[\mathbin{\overset{\mathrm L}{\otimes}}_S:D^-(f_1)_{\mathrm{coh}}\times D^-(f_2)_{\mathrm{coh}}\longrightarrow D^-(f)_{\mathrm{coh}}\qquad(\text{见 1.2}).\]
特别地，若 $f_1,f_2$ \Term{pseudo-coherent}{伪凝聚}（1.2），则 $f$ 也如此。
\end{SGAresult}
\begin{SGAresult}{推论}{1.10}
设有\Term{cartesian-square}{笛卡尔方块}
\[
\begin{tikzcd}X\arrow[d,"f"']&X'\arrow[l,"g"']\arrow[d,"f'"]\\Y&Y'\arrow[l,"h"]\end{tikzcd}
\]
其中 $f$ \Term{locally-finite-type}{局部有限型}。则：

a) 假定 $X,Y'$ 在 $Y$ 上\Term{tor-independent}{Tor 无关}（1.5）。对 $n\in\mathbf Z\cup\{-\infty\}$，\Term{functor}{函子} $\mathrm Lg^*:D^-(X)\to D^-(X')$ 诱导\Term{functor}{函子}
\[\mathrm Lg^*:D^-(f)_{n\text{-coh}}\longrightarrow D^-(f')_{n\text{-coh}}.\]
特别地，若 $f$ 为 $n$-\Term{pseudo-coherent}{伪凝聚}（1.2），$f'$ 也如此。

b) 假定 $h$ \Term{faithfully-flat}{忠实平坦}。设 $n\in\mathbf Z\cup\{-\infty\}$，$E\in\operatorname{Ob}D^-(X)$ 是具有\Term{quasi-coherent}{拟凝聚}\Term{cohomology}{上同调}的\Term{complex}{复形}。则 $E$ \Term{relative-pseudo-coherent}{相对于 $f$ 为 $n$-伪凝聚}，当且仅当 $g^*E$ \Term{relative-pseudo-coherent}{相对于 $f'$ 为 $n$-伪凝聚}。特别地，$f$ 为 $n$-\Term{pseudo-coherent}{伪凝聚}，当且仅当 $f'$ 如此。
\end{SGAresult}
\begin{proof}
考虑到 1.2.1 b)，1.9、1.10 中所有断言均直接来自 1.8，惟 1.10 b) 的充分性除外。为证这一点，因问题在 $X$ 上是\Term{local-scheme}{局部}的，可假定 $f$ \Term{resolution}{分解}为\Term{closed-immersion}{闭浸入}与\Term{smooth}{光滑态射}的复合。利用 1.1.3，遂归结为 $f$（相应地 $f'$）为 $X$（相应地 $X'$）上的恒等\Term{morphism}{态射}的情形。再作标准的逐步归约（参见第 I 讲 2.16.2），利用 $X$ 上的\Term{quasi-coherent}{拟凝聚}\Term{sheaf}{层} $L$ 为\Term{finite-type}{有限型}当且仅当 $g^*L$ 为\Term{finite-type}{有限型}，即得断言。
\end{proof}
\SourcePage{244}
\begin{SGAresult}{命题}{1.11}
设
\[
\begin{tikzcd}X\arrow[r,"f"]\arrow[dr,"h"']&Y\arrow[d,"g"]\\&Z\end{tikzcd}
\]
为\Term{locally-finite-type}{局部有限型}\Term{morphism}{态射}的\Term{commutative-diagram}{交换三角形}；设 $E\in\operatorname{Ob}D(X)$，$F\in\operatorname{Ob}D(Y)$，$[a,b],[c,d]$ 为 $\mathbf Z$ 中的区间。

a) 若 $E$ \Term{relative-pseudo-coherent}{相对于 $f$ 为 $a$-伪凝聚}且在\Term{degree}{次数} $>b$ \Term{acyclic}{无上同调}，$F$ \Term{relative-pseudo-coherent}{相对于 $g$ 为 $c$-伪凝聚}且在\Term{degree}{次数} $>d$ \Term{acyclic}{无上同调}，则 $E\mathbin{\overset{\mathrm L}{\otimes}}_YF$（见 1.6）\Term{relative-pseudo-coherent}{相对于 $h$ 为 $m$-伪凝聚}且在\Term{degree}{次数} $>n$ \Term{acyclic}{无上同调}，其中 $m=\sup(a+d,b+c)$，$n=b+d$。

b) 若 $E$ \Term{relative-pseudo-coherent}{相对于 $h$ 为 $a$-伪凝聚}且在\Term{degree}{次数} $>b$ \Term{acyclic}{无上同调}，而 $g$ 为 $(a+1-b)$-\Term{pseudo-coherent}{伪凝聚}，则 $E$ \Term{relative-pseudo-coherent}{相对于 $f$ 为 $a$-伪凝聚}。
\end{SGAresult}
\begin{proof}
问题在 $X$ 上是\Term{local-scheme}{局部}的，可假定有如下\Term{commutative-diagram}{交换图}，其中方块为\Term{cartesian-square}{笛卡尔方块}：
\[
\begin{tikzcd}[column sep=large,row sep=large]
&\mathbf E_Y^r\arrow[r,"j"]\arrow[d,"p"]&\mathbf E_Z^{r+s}\arrow[d]\\
X\arrow[ur,"i"]\arrow[r,"f"']&Y\arrow[r,"k"]\arrow[dr,"g"']&\mathbf E_Z^s\arrow[d,"q"]\\&&Z
\end{tikzcd}
\]
其中 $i,j,k$ 为\Term{closed-immersion}{闭浸入}，$p,q$ 为标准\Term{affine-space}{仿射空间}的\Term{canonical-projection}{典范投影}。若 $F$ \Term{relative-pseudo-coherent}{相对于 $Z$ 为 $c$-伪凝聚}，则由 1.1 (iii)，$k_*F$ 是 $c$-\Term{pseudo-coherent}{伪凝聚}的，因而由 1.1.3，$j_*p^*F$ 也如此。同样，若 $g$ 为 $x$-\Term{pseudo-coherent}{伪凝聚}，则 $j_*\mathcal O_{\mathbf E_Y^r}$ 为 $x$-\Term{pseudo-coherent}{伪凝聚}。根据\Term{relative-pseudo-coherent}{相对伪凝聚性}的定义（1.2），证明 a)、b) 遂归结为 $f,g$ 均为\Term{closed-immersion}{闭浸入}的情形。
\SourcePage{245}
\textbf{a) 的证明。} 由 1.1 (iii)，须证明：若 $f_*E$ 为 $a$-\Term{pseudo-coherent}{伪凝聚}且在\Term{degree}{次数} $>b$ \Term{acyclic}{无上同调}，$g_*F$ 为 $c$-\Term{pseudo-coherent}{伪凝聚}且在\Term{degree}{次数} $>d$ \Term{acyclic}{无上同调}，则 $h_*(E\mathbin{\overset{\mathrm L}{\otimes}}_YF)$ 为 $m$-\Term{pseudo-coherent}{伪凝聚}且在\Term{degree}{次数} $>n$ \Term{acyclic}{无上同调}。有\Term{canonical}{典范}\Term{isomorphism}{同构}
\[f_*E\mathbin{\overset{\mathrm L}{\otimes}}F\xrightarrow{\sim}f_*(E\mathbin{\overset{\mathrm L}{\otimes}}_YF),\]
这是 1.8 证明中\Term{isomorphism}{同构} (*) 的特例。\Term{functor}{函子} $g_*(-\mathbin{\overset{\mathrm L}{\otimes}}F)$ 把\Term{finite-type}{有限型}\Term{locally-free}{局部自由}\Term{sheaf}{层}变成 $c$-\Term{pseudo-coherent}{伪凝聚}\Term{complex}{复形}，并把在\Term{degree}{次数} $>0$ \Term{acyclic}{无上同调}的\Term{complex}{复形}变成在\Term{degree}{次数} $>d$ \Term{acyclic}{无上同调}的\Term{complex}{复形}，所以断言来自第 I 讲 2.13。

\textbf{b) 的证明。} 须证明：若 $h_*E$ 是 $a$-\Term{pseudo-coherent}{伪凝聚}且在\Term{degree}{次数} $>b$ \Term{acyclic}{无上同调}，$g_*\mathcal O_Y$ 是 $(a+1-b)$-\Term{pseudo-coherent}{伪凝聚}的，则 $f_*E$ 为 $a$-\Term{pseudo-coherent}{伪凝聚}。既然断言只涉及 $f_*E$，可假定 $X=Y$、$f=\operatorname{Id}_X$。像 1.1.1 的证明那样，作\Term{descent}{下降}归纳，证明 $E$ 在\Term{degree}{次数} $\ge a$ \Term{pseudo-coherent}{伪凝聚}。假定已经证明 $E$ 为 $(a+1)$-\Term{pseudo-coherent}{伪凝聚}，则\Term{localization}{局部化}后由第 I 讲 2.10 有\Term{distinguished-triangle}{可区别三角形}
\[E'\longrightarrow E\longrightarrow E''\longrightarrow E'[1],\]
其中当 $i\in[a+1,b]$ 时 $E'{}^i$ \Term{locally-free}{局部自由}，当 $i\notin[a+1,b]$ 时 $E'{}^i=0$，而 $i>a$ 时 $E''{}^i=0$。$g_*\mathcal O_Y$ 为 $(a+1-b)$-\Term{pseudo-coherent}{伪凝聚}的假设，由第 I 讲 2.11 a) 蕴含 $g_*E'$ 为 $(a+1)$-\Term{pseudo-coherent}{伪凝聚}。由于 $g_*E$ 是 $a$-\Term{pseudo-coherent}{伪凝聚}的，第 I 讲 2.6 给出 $g_*E''$ 也是 $a$-\Term{pseudo-coherent}{伪凝聚}的。再由第 I 讲 2.10 b)，并利用一个 $\mathcal O_Y$-模 $M$ 为\Term{finite-type}{有限型}当且仅当 $g_*M$ 为\Term{finite-type}{有限型}，即可完成证明。
\end{proof}
由 1.11 直接得到下列推论。
\SourcePage{246}
\begin{SGAresult}{推论}{1.12}
在 1.11 的情形下，若 $g$ \Term{pseudo-coherent}{伪凝聚}，则对 $r\in\mathbf Z\cup\{-\infty\}$，有
\[D^-(f)_{r\text{-coh}}=D^-(h)_{r\text{-coh}}.\]
特别地（令 $X=Y$），在 $Y$ 上，“绝对意义的”$r$-\Term{pseudo-coherent}{伪凝聚性}与\Term{relative-pseudo-coherent}{相对于 $g$ 的 $r$-伪凝聚性}一致。
\end{SGAresult}
\begin{SGAresult}{推论}{1.13}
在 1.11 的情形下，若 $f$ \Term{pseudo-coherent}{伪凝聚}，则\Term{functor}{函子} $\mathcal O_X\mathbin{\overset{\mathrm L}{\otimes}}_Y- =\mathrm Lf^*:D^-(Y)\to D^-(X)$ 对每个 $r\in\mathbf Z\cup\{-\infty\}$ 诱导\Term{functor}{函子}
\[\mathrm Lf^*:D^-(g)_{r\text{-coh}}\longrightarrow D^-(h)_{r\text{-coh}}.\]
\end{SGAresult}
\begin{SGAresult}{推论}{1.14}
在 1.11 的情形下，若 $f$ 是 $r$-\Term{pseudo-coherent}{伪凝聚}的，而 $g$ 是 $s$-\Term{pseudo-coherent}{伪凝聚}的（$r,s\in\mathbf Z\cup\{-\infty\}$），则 $h$ 是 $\sup(r,s)$-\Term{pseudo-coherent}{伪凝聚}的。若 $h$ 是 $r$-\Term{pseudo-coherent}{伪凝聚}的，$g$ 是 $(r+1)$-\Term{pseudo-coherent}{伪凝聚}的，则 $f$ 是 $r$-\Term{pseudo-coherent}{伪凝聚}的。
\end{SGAresult}
\textbf{注记 1.14.1。} 特别地，当 $r=s=-1$ 时，便重新得到 EGA IV 1.4.3 的断言 (ii)、(v)（参见 1.2.1 c)）。
\SGAsection{2}{有限性定理}
\begin{SGAresult}{猜想}{2.1}
设 $S$ 为\Term{scheme}{概形}，$f:X\to Y$ 为在 $S$ 上\Term{locally-finite-type}{局部有限型}的\Term{scheme}{概形}之间的\Term{proper}{固有} $S$-\Term{morphism}{态射}，$E\in\operatorname{Ob}D^+(X)$。若 $E$ \Term{relative-pseudo-coherent}{相对于 $S$ 伪凝聚}（1.2），则 $\mathrm Rf_*E$ 也如此。
\end{SGAresult}
\begin{SGAresult}{\Term{theorem}{定理}}{2.2}
上述猜想在下列情形成立：
\begin{enumerate}\item[(i)] $S$ \Term{locally-noetherian}{局部 Noether}；\item[(ii)] $f$ \Term{projective}{射影}。\end{enumerate}
\end{SGAresult}
\SourcePage{247}
\textbf{情形 (i) 的证明。} 因 $S$ \Term{locally-noetherian}{局部 Noether}，$X$（相应地 $Y$）到 $S$ 的\Term{projection}{投影}是\Term{pseudo-coherent}{伪凝聚}\Term{morphism}{态射}（1.3），而 $X$（相应地 $Y$）上的\Term{pseudo-coherent}{伪凝聚性}与\Term{relative-pseudo-coherent}{相对于 $S$ 的伪凝聚性}一致（1.12）。另一方面，$X$（相应地 $Y$）\Term{locally-noetherian}{局部 Noether}，故 $\mathcal O_X$（相应地 $\mathcal O_Y$）\Term{coherent}{凝聚}；因此由第 I 讲 3.5，$D(X)$（相应地 $D(Y)$）的\Term{object}{对象} $M$ \Term{pseudo-coherent}{伪凝聚}，当且仅当 $H^i(M)$ \Term{coherent}{凝聚}且 $M$ \Term{local-scheme}{局部}属于 $D^-$。所以本\Term{theorem}{定理}来自下述结果。
\begin{SGAresult}{\Term{theorem}{定理}}{2.2.1}
设 $f:X\to Y$ 是\Term{locally-noetherian}{局部 Noether}\Term{scheme}{概形}的\Term{proper}{固有态射}，$E\in\operatorname{Ob}D^+(X)$。若 $H^i(E)$ \Term{coherent}{凝聚}，则 $R^if_*E$ 也如此。若再假定 $E$ \Term{local-scheme}{局部}属于 $D^-$，则 $\mathrm Rf_*E$ 也如此。
\end{SGAresult}
\begin{proof}
第一断言来自\Term{spectral-sequence}{谱序列}
\[R^pf_*\bigl(H^q(E)\bigr)\Longrightarrow R^*f_*E\tag{*}\]
以及\Term{finiteness-theorem}{有限性定理}（EGA III 3.2.1）。对于第二断言，可假定 $Y$ \Term{affine}{仿射}；此时 $X$ \Term{quasi-compact}{拟紧}且\Term{separated}{分离}，设它可由 $r$ 个\Term{affine}{仿射}\Term{open-subset}{开集}\Term{covering}{覆盖}。已知在此情形下，对每个\Term{quasi-coherent}{拟凝聚} $\mathcal O_X$-模 $M$ 及每个 $i>r$，有 $R^if_*M=0$（EGA III 1.4.12）。若 $E$ \Term{local-scheme}{局部}属于 $D^-$，因 $X$ \Term{quasi-compact}{拟紧}，便有 $E\in D^-$；设 $H^i(E)=0$（$i>N$）。由\Term{spectral-sequence}{谱序列} (*) 得 $R^if_*E=0$（$i>N+r$），即得所证。
\end{proof}
\textbf{情形 (ii) 的证明。} 问题在 $Y$ 上是\Term{local-scheme}{局部}的，可假定 $S,Y$ 均\Term{affine}{仿射}。于是有\Term{commutative-diagram}{交换图}，其中方块为\Term{cartesian-square}{笛卡尔方块}：
\[
\begin{tikzcd}[column sep=large,row sep=large]&\mathbf P_Y^r\arrow[r,"j"]\arrow[d,"p"]&\mathbf P_{Y'}^r\arrow[d,"p'"]\\X\arrow[ur,"i"]\arrow[r,"f"']&Y\arrow[r,"j'"]\arrow[dr,"g"']&Y'\arrow[d,"g'"]\\&&S\end{tikzcd}
\]
\SourcePage{248}
其中 $i,j,j'$ 为\Term{closed-immersion}{闭浸入}，$g'$ 是\Term{affine}{仿射}\Term{smooth}{光滑态射}，$p,p'$ 为标准\Term{projective-space}{射影空间}的\Term{projection}{投影}。假定 $E$ \Term{relative-pseudo-coherent}{相对于 $S$ 伪凝聚}。由 1.1 (iii)，这意味着 $(ji)_*E$ \Term{pseudo-coherent}{伪凝聚}。须证明 $\mathrm Rf_*E$ \Term{relative-pseudo-coherent}{相对于 $S$ 伪凝聚}，即由 1.1 (iii)，须证明 $j'_*\mathrm Rf_*E$ \Term{pseudo-coherent}{伪凝聚}。但有
\[j'_*\mathrm Rf_*E\xrightarrow{\sim}\mathrm Rp'_*(ji)_*E,\]
所以断言来自下述更精确的结果。
\begin{SGAresult}{\Term{theorem}{定理}}{2.2.2}
设 $Y$ 为\Term{affine}{仿射}\Term{scheme}{概形}，$f:X=\mathbf P_Y^r\to Y$ 为标准\Term{projective-space}{射影空间}的\Term{projection}{投影}，$E\in\operatorname{Ob}D^b(X)_{\mathrm{coh}}$，$a\in\mathbf Z$。则 $\mathrm Rf_*E$ \Term{pseudo-coherent}{伪凝聚}。此外，若 $E$ 在\Term{degree}{次数} $>a$ \Term{acyclic}{无上同调}，则 $\mathrm Rf_*E$ 在\Term{degree}{次数} $>a+r$ \Term{acyclic}{无上同调}，且存在整数 $N$，使对所有 $n\ge N$，$\mathrm Rf_*(E(n))$ 在\Term{degree}{次数} $>a$ \Term{acyclic}{无上同调}，其中 $E(n)=E\otimes\mathcal O_X(n)$。
\end{SGAresult}
\begin{proof}
因 $Y$ \Term{affine}{仿射}，$X$ \Term{quasi-compact}{拟紧}且\Term{separated}{分离}，故由第 II 讲 2.2.9 a)，$E$ 在 $D(X)$ 中\Term{isomorphism}{同构}于一个有上界的\Term{complex}{复形} $F$，其每个 $F^i$ 属于 $\mathsf L(X)$，即（同上引用处）为形如 $\mathcal O_X(-n)$（$n>0$）的\Term{sheaf}{层}的\Term{finite}{有限}\Term{direct-sum}{直和}。此外，若 $E$ 在\Term{degree}{次数} $>a$ \Term{acyclic}{无上同调}，则可取 $F$ 使 $F^i=0$（$i>a$）：结合第 I 讲 2.10 a) 与第 II 讲 2.2.8 a)，这在\Term{form}{形式}上来自前述断言。用 $F$ 代替 $E$，可假定每个 $E^i\in\operatorname{Ob}\mathsf L(X)$，且 $E^i=0$（$i>a$）。因为 $X$ 可由 $r+1$ 个\Term{affine}{仿射}\Term{open-subset}{开集}\Term{covering}{覆盖}，由 EGA III 1.4.10 和超\Term{cohomology}{上同调}\Term{spectral-sequence}{谱序列}，有 $R^qf_*E=0$（$q>a+r$）。取整数 $t\le a$，令
\[E'=(0\longrightarrow E^t\longrightarrow\cdots\longrightarrow E^a\longrightarrow0).\]
有\Term{exact-sequence}{正合列} $0\to E'\to E\to E''\to0$，
\SourcePage{249}
其中 $E''$ 满足与 $E$ 相同的条件，只须将 $a$ 换成 $t-1$。施行 $\mathrm Rf_*$，得到\Term{distinguished-triangle}{可区别三角形}
\[\mathrm Rf_*E'\longrightarrow\mathrm Rf_*E\longrightarrow\mathrm Rf_*E''\longrightarrow(\mathrm Rf_*E')[1].\tag{*}\]
$t$ 可以任意小，因此对于第一断言，只须证明 $\mathrm Rf_*E$ 为 $(t+r)$-\Term{pseudo-coherent}{伪凝聚}。$\mathrm Rf_*E''$ 在\Term{degree}{次数} $\ge t+r$ \Term{acyclic}{无上同调}，故为 $(t+r)$-\Term{pseudo-coherent}{伪凝聚}。由第 I 讲 2.6，只要证明 $\mathrm Rf_*E'$ \Term{perfect-complex}{完美}即可。为此，逐步归约（第 I 讲 4.10），可化为 $E'$ 仅在\Term{degree}{次数}零非零且等于 $\mathcal O_X(-n)$（$n>0$）的情形。由 EGA III 2.1.16，$R^rf_*\mathcal O_X(-n)$ 为\Term{finite-type}{有限型}\Term{locally-free}{局部自由}，且
\[\mathrm Rf_*\mathcal O_X(-n)\xrightarrow{\sim}R^rf_*\mathcal O_X(-n)[-r].\]
所以 $\mathrm Rf_*\mathcal O_X(-n)$ \Term{perfect-complex}{完美}，由此得到 2.2.2 的第一断言。已知 $\mathrm Rf_*E$ 在\Term{degree}{次数} $>a+r$ \Term{acyclic}{无上同调}。尚须证明最后一条断言。当 $t<-r+a$ 时，$\mathrm Rf_*E''$ 在\Term{degree}{次数} $>a$ \Term{acyclic}{无上同调}，故结合三角形 (*)，可用 $E'$ 代替 $E$ 来证明。因为 $E'$ \Term{degree}{次数}\Term{bounded}{有界}，各\Term{set-component}{分量}属于 $\mathsf L(X)$，存在整数 $N$，使当 $n\ge N$ 时，$E'(n)$ 的各\Term{set-component}{分量}均为形如 $\mathcal O_X(p)$（$p\ge0$）的\Term{sheaf}{层}的\Term{finite}{有限}\Term{direct-sum}{直和}。由 EGA III 2.1.16 和超\Term{cohomology}{上同调}\Term{spectral-sequence}{谱序列}，对所有 $n\ge N$，$\mathrm Rf_*(E'(n))$ 在\Term{degree}{次数} $>a$ \Term{acyclic}{无上同调}。这完成了 2.2.2 的证明，因而也完成了 2.2 的证明。
\end{proof}
\Term{theorem}{定理} 2.2 (ii) 有一个“经典”的补充，列于下文以备后用。
\SourcePage{250}
\begin{SGAresult}{推论}{2.3（EGA III 2.2.1）}
设 $S$ 为\Term{quasi-compact}{拟紧}\Term{scheme}{概形}，$f:X\to Y$ 为 $S$ 上\Term{finite-type}{有限型}\Term{scheme}{概形}之间的\Term{projective}{射影} $S$-\Term{morphism}{态射}；设 $L=i^*\mathcal O_P(1)$ 为\Term{embedding}{嵌入} $i:X\to P=\mathbf P(F)$ 所定义的\Term{invertible-sheaf}{可逆层}，其中 $F$ 为\Term{finite-type}{有限型}\Term{quasi-coherent}{拟凝聚} $\mathcal O_Y$-模。对 $M\in\operatorname{Ob}D(X)$、$n\in\mathbf Z$，令 $M(n)=M\otimes L^{\otimes n}$。若 $E\in\operatorname{Ob}D^+(X)$ \Term{relative-pseudo-coherent}{相对于 $S$ 伪凝聚}，且在\Term{degree}{次数} $>a$ \Term{acyclic}{无上同调}（$a\in\mathbf Z$），则存在整数 $N$，使对所有 $n\ge N$，$\mathrm Rf_*(E(n))$ 满足相同条件。
\end{SGAresult}
\begin{proof}
像 2.2 (ii) 的证明一样，直接归结为 2.2.2 的情形：$Y=S$，$f:X=\mathbf P_Y^r\to Y$ 为\Term{canonical-projection}{典范投影}，$L=\mathcal O_X(1)$ 为\Term{canonical}{典范}\Term{sheaf}{层}，$E\in\operatorname{Ob}D^b(X)$ 为\Term{pseudo-coherent}{伪凝聚}且在\Term{degree}{次数} $>a$ \Term{acyclic}{无上同调}的\Term{complex}{复形}；于是 2.2.2 给出结论。
\end{proof}
\begin{SGAresult}{推论}{2.4}
设 $S$ 为\Term{quasi-compact}{拟紧}\Term{scheme}{概形}，$f:X\to Y$ 为 $S$ 上\Term{finite-type}{有限型}\Term{scheme}{概形}之间的\Term{proper}{固有} $S$-\Term{morphism}{态射}。若 $S$ 为\Term{noetherian}{Noether}，或 $f$ \Term{projective}{射影}，则\Term{functor}{函子} $\mathrm Rf_*:D^+(X)\to D^+(Y)$ 诱导\Term{functor}{函子}
\[\mathrm Rf_*:D^b(X)_{S\text{-coh}}\longrightarrow D^b(Y)_{S\text{-coh}}.\]
\end{SGAresult}
\begin{proof}
由 2.2 以及以下事实立即得到：此时存在整数 $N$，使对每个\Term{quasi-coherent}{拟凝聚} $\mathcal O_X$-模 $E$ 及每个 $i>N$，有 $R^if_*E=0$（EGA III 1.4.12、IV 1.7.21）。
\end{proof}
\begin{SGAresult}{推论}{2.5}
设 $f:X\to Y$ 是\Term{proper}{固有}且\Term{pseudo-coherent}{伪凝聚}的\Term{morphism}{态射}（1.2），$E\in\operatorname{Ob}D^+(X)$。假定 $Y$ \Term{locally-noetherian}{局部 Noether}，或 $f$ \Term{projective}{射影}。若 $E$（在绝对意义下）\Term{pseudo-coherent}{伪凝聚}，则 $\mathrm Rf_*E$ 也如此。若 $Y$ \Term{quasi-compact}{拟紧}，\Term{functor}{函子} $\mathrm Rf_*$ 诱导\Term{functor}{函子}
\[\mathrm Rf_*:D^b(X)_{\mathrm{coh}}\longrightarrow D^b(Y)_{\mathrm{coh}}.\]
\end{SGAresult}
\SourcePage{251}
\begin{proof}这是 2.2、2.5 和 1.12 的推论。\end{proof}
\textbf{评注 2.6。} 要在\Term{generic}{一般}情形下证明猜想 2.1，无疑需要采用新的技术。我们希望这些技术有一天也有助于证明\Term{analytic-geometry}{解析几何}中的类似猜想：对于一个以具有适当\Term{topology}{拓扑}的 $\mathbf C$-\Term{algebra}{代数}为结构环的\Term{locally-ringed-topos}{局部赋环拓扑斯}之上的\Term{analytic-space}{解析空间}的\Term{proper}{固有态射}，有相应的结论。到目前为止，此猜想仅在底为\Term{complex-analytic-space}{复解析空间}时得到了“证明”（Grauert \Term{theorem}{定理} [1]）。\footnote{（1971 年 4 月增补）参见 R. Kiehl 与 J. L. Verdier 正在准备的一篇文章；其中以另一种更简单的技术，在更\Term{generic}{一般}的假设下证明了\Term{finiteness-theorem}{有限性定理} 1。}
\SGAsection{3}{相对有限 Tor 维数}
\begin{SGAresult}{定义}{3.1}
设 $f:(X,\mathcal O_X)\to(Y,\mathcal O_Y)$ 为\Term{ringed-topos}{赋环拓扑斯}的\Term{morphism}{态射}，$E\in\operatorname{Ob}D^-(X)$，$[m,n]$ 为 $\mathbf Z$ 中的区间。若把 $E$ 看作 $f^{-1}\mathcal O_Y$-模\Term{complex}{复形}时，它的\Term{flat-amplitude}{平坦振幅}包含于 $[m,n]$（第 I 讲 5.2），则称 $E$ \Term{relative-flat-amplitude}{相对于 $f$ 的平坦振幅}（无歧义时也说相对于 $Y$）包含于 $[m,n]$，记作
\[\operatorname{tor.amp}_f(E)\subset[m,n]\qquad\bigl(\text{或 }\operatorname{tor.amp}_Y(E)\subset[m,n]\bigr).\]
若 $\mathcal O_X$ \Term{relative-flat-amplitude}{相对于 $f$ 的平坦振幅}包含于 $[-a,0]$（$a\ge0$），则称 $f$ 的\Term{flat-amplitude}{平坦振幅}包含于 $[-a,0]$，或称它的\Term{tor-dimension}{Tor 维数} $\le a$。若存在区间 $[m,n]$ 使 $\operatorname{tor.amp}_f(E)\subset[m,n]$，则称 $E$ \Term{relative-flat-amplitude}{相对于 $f$ 有限平坦振幅}。
\end{SGAresult}
\textbf{记号 3.2。} 记 $D(f)_{\mathrm{torf}}$（或 $D(X)_{f\text{-torf}}$，无混淆时也记作 $D(X)_{Y\text{-torf}}$）为由\Term{relative-flat-amplitude}{相对于 $f$ 有限平坦振幅}的\Term{complex}{复形}组成的 $D^-(X)$ 的\Term{full-subcategory}{满子范畴}。由第 I 讲 5.3，它是 $D^b(X)$ 的\Term{triangulated-subcategory}{三角子范畴}。
\SourcePage{252}
\begin{SGAresult}{命题}{3.3}
在 3.1 的条件下，假定 $X$ 有\Term{enough-points}{足够多的点}（SGA 4 第 XVII 讲；第 IV 讲 6.4）。则下列条件\Term{equivalence}{等价}：
\begin{enumerate}
\item[(i)] $\operatorname{tor.amp}_f(E)\subset[m,n]$；
\item[(ii)] 对 $X$ 的每个点 $x$，有 $\operatorname{tor.amp}_y(E_x)\subset[m,n]$，其中 $y$ 为 $x$ 在 $f$ 下的像点，赋予环 $\mathcal O_{Y,y}$；
\item[(iii)] 对 $Y$ 的每个\Term{object}{对象} $V$ 和每个 $\mathcal O_V$-模 $M$，有
\[H^i\bigl((E|f^{-1}V)\mathbin{\overset{\mathrm L}{\otimes}}_VM\bigr)=0\qquad(i\notin[m,n]),\]
其中 $(E|f^{-1}V)\mathbin{\overset{\mathrm L}{\otimes}}_VM=(E|f^{-1}V)\mathbin{\overset{\mathrm L}{\otimes}}\mathrm Lf^*M$；
\item[(iv)] 与 (iii) 相同，但只要求 $M$ 是某个\Term{homomorphism}{同态} $\mathcal O_V^q\to\mathcal O_V^p$ 的\Term{cokernel}{余核}，其中 $p,q\in\mathbf N$。
\end{enumerate}
\end{SGAresult}
\begin{proof}
(i)$\Longleftrightarrow$(ii) 来自第 I 讲 5.6.2，而 (ii)$\Longrightarrow$(iii)$\Longrightarrow$(iv) 显然。证明 (iv)$\Longrightarrow$(ii)。设 $x$ 为 $X$ 的一点，$y=f(x)$；须证明对每个 $\mathcal O_{Y,y}$-模 $M$ 及每个 $i\notin[m,n]$，有
\[H^i\bigl(E_x\mathbin{\overset{\mathrm L}{\otimes}}_{\mathcal O_{Y,y}}M\bigr)=0.\tag{*}\]
每个 $\mathcal O_{Y,y}$-模都是\Term{finite-presentation}{有限表示}模的\Term{filtered-colimit}{滤过归纳极限}，故只须对\Term{finite-presentation}{有限表示}的 $M$ 验证上式。若 $M$ 为\Term{finite-presentation}{有限表示}的 $\mathcal O_{Y,y}$-模，采用与普通\Term{ringed-space}{赋环空间}相同的论证（EGA $0_{\mathrm I}$ 5.3.9），可得 $Y$ 的一个带 $y$-点的\Term{object}{对象} $V$，以及一个为\Term{homomorphism}{同态} $\mathcal O_V^q\to\mathcal O_V^p$（$p,q\in\mathbf N$）之\Term{cokernel}{余核}的 $\mathcal O_V$-模 $P$，使 $P_y\simeq M$。此时
\[\bigl((E|f^{-1}V)\mathbin{\overset{\mathrm L}{\otimes}}_VP\bigr)_x\xrightarrow{\sim}E_x\mathbin{\overset{\mathrm L}{\otimes}}_{\mathcal O_{Y,y}}M,\]
\SourcePage{253}
因为\Term{tensor-product}{张量积}与取\Term{stalk}{茎}可交换；假设 (iv) 便蕴含 (*)，证毕。
\end{proof}
\textbf{注记 3.3.1。} 撰稿人不知道，去掉 $X$ 有\Term{enough-points}{足够多的点}这一假设后，(i)$\Longleftrightarrow$(iii)$\Longleftrightarrow$(iv) 是否仍然成立。
\begin{SGAresult}{推论}{3.4}
设
\[
\begin{tikzcd}X\arrow[r,"f"]\arrow[dr,"h"']&Y\arrow[d,"g"]\\&Z\end{tikzcd}
\]
为\Term{ringed-topos}{赋环拓扑斯}的\Term{commutative-diagram}{交换三角形}，假定 $X$ 有\Term{enough-points}{足够多的点}。设 $E\in\operatorname{Ob}D(X)$，$F\in\operatorname{Ob}D(Y)$，$[a,b],[c,d]$ 为 $\mathbf Z$ 中的区间。则
\[\left.\begin{gathered}\operatorname{tor.amp}_f(E)\subset[a,b]\\\operatorname{tor.amp}_g(F)\subset[c,d]\end{gathered}\right\}
\Longrightarrow\operatorname{tor.amp}_h(E\mathbin{\overset{\mathrm L}{\otimes}}_YF)\subset[a+c,b+d].\]
\end{SGAresult}
\begin{proof}
由 3.3，须证明左侧的关系蕴含：对每个 $i\notin[a+c,b+d]$ 和每个 $\mathcal O_{Z'}$-模 $M$，其中 $Z'$ 为 $Z$ 的\Term{object}{对象}，有
\[H^i\bigl(((E\mathbin{\overset{\mathrm L}{\otimes}}_YF)|h^{-1}Z')\mathbin{\overset{\mathrm L}{\otimes}}_{Z'}M\bigr)=0.\]
令 $Y'=g^{-1}Z'$、$X'=f^{-1}Y'=h^{-1}Z'$，则
\[\bigl((E\mathbin{\overset{\mathrm L}{\otimes}}_YF)|X'\bigr)\mathbin{\overset{\mathrm L}{\otimes}}_{Z'}M
\simeq(E|X')\mathbin{\overset{\mathrm L}{\otimes}}_{Y'}\bigl((F|Y')\mathbin{\overset{\mathrm L}{\otimes}}_{Z'}M\bigr),\]
所以断言立即由假设得到。
\end{proof}
特别地：
\begin{SGAresult}{推论}{3.4.1}
在 3.4 的假设下，

a) 若 $g$ 的\Term{tor-dimension}{Tor 维数} $\le n$，则
\[\operatorname{tor.amp}_f(E)\subset[a,b]\Longrightarrow\operatorname{tor.amp}_h(E)\subset[a-n,b];\]
\SourcePage{254}
b) 若 $f$ 的\Term{tor-dimension}{Tor 维数} $\le m$，则
\[\operatorname{tor.amp}_g(F)\subset[c,d]\Longrightarrow\operatorname{amp}_h(\mathrm Lf^*F)\subset[c-m,d].\]
\end{SGAresult}
\textbf{注记 3.4.2。} 设 $f:X\to Y$ 为\Term{finite}{有限}\Term{tor-dimension}{Tor 维数}的\Term{morphism}{态射}，$E\in\operatorname{Ob}D(X)$。由 3.4.1 a)，若 $E$（在绝对意义下）有\Term{finite}{有限}\Term{flat-amplitude}{平坦振幅}，则 $E$ \Term{relative-flat-amplitude}{相对于 $f$ 有限平坦振幅}。但逆命题不成立，即使 $f$ \Term{flat}{平坦}也是如此；域上的一个\Term{scheme}{概形}的\Term{singular-point}{奇点}就提供了例子。
\begin{SGAresult}{命题}{3.5}
设 $S$ 为\Term{scheme}{概形}；对 $i=1,2$，设 $f_i:X_i\to Y_i$ 为 $S$-\Term{morphism}{态射}，$E_i\in\operatorname{Ob}D^-(X_i)$，$[a_i,b_i]$ 为 $\mathbf Z$ 中的区间。假定 $X_1,X_2$ 以及 $Y_1,Y_2$ 在 $S$ 上\Term{tor-independent}{Tor 无关}（1.5）。令
\[a=a_1+a_2,\quad b=b_1+b_2,\quad f=f_1\times_Sf_2:X\longrightarrow Y.\]
a) 若 $\operatorname{tor.amp}_{f_1}(E_1)\subset[a_1,b_1]$，且 $E_2$ 在 $[a_2,b_2]$ 之外\Term{acyclic}{无上同调}，则 $E_1\mathbin{\overset{\mathrm L}{\otimes}}_SE_2$（见 1.6）在 $[a,b]$ 之外\Term{acyclic}{无上同调}。

b) 若 $i=1,2$ 均满足 $\operatorname{tor.amp}_{f_i}(E_i)\subset[a_i,b_i]$，则
\[\operatorname{tor.amp}_f(E_1\mathbin{\overset{\mathrm L}{\otimes}}_SE_2)\subset[a,b].\]
\end{SGAresult}
\begin{proof}
由 3.3，只须在“\Term{stalk}{茎}上”验证。设 $s$ 为 $S$ 的一点，$x$ 为 $X$ 中位于 $s$ 之上的一点，$y=f(x)$，$x_i$（相应地 $y_i$）为 $x$（相应地 $y$）在 $X_i$（相应地 $Y_i$）上的\Term{canonical-projection}{典范投影}。令
\[
\begin{gathered}
(E_i)_{x_i}=M_i,\quad\mathcal O_{X_i,x_i}=A_i,\quad\mathcal O_{Y_i,y_i}=B_i,\quad\mathcal O_{S,s}=C,\\
\mathcal O_{X,x}=A=A_1\otimes_CA_2,\qquad \mathcal O_{Y,y}=B=B_1\otimes_CB_2.
\end{gathered}
\]
有
\[(E_1\mathbin{\overset{\mathrm L}{\otimes}}_SE_2)_x\simeq M_1\mathbin{\overset{\mathrm L}{\otimes}}_{A_1}A\mathbin{\overset{\mathrm L}{\otimes}}_{A_2}M_2.\tag{*}\]
另一方面，\Term{tor-independent}{Tor 无关}的假设蕴含
\[A\simeq A_1\mathbin{\overset{\mathrm L}{\otimes}}_{B_1}B\mathbin{\overset{\mathrm L}{\otimes}}_{B_2}A_2.\]
所以 (*) 的右侧在 $B$-模的\Term{derived-category}{导出范畴} $D(B)$ 中\Term{canonical}{典范}地\Term{isomorphism}{同构}于
\SourcePage{255}
$M_1\mathbin{\overset{\mathrm L}{\otimes}}_{B_1}B\mathbin{\overset{\mathrm L}{\otimes}}_{B_2}M_2$。关系 $\operatorname{tor.amp}_{f_i}(E_i)\subset[a_i,b_i]$（$i=1,2$）蕴含 $\operatorname{tor.amp}_{B_i}(M_i)\subset[a_i,b_i]$，故由第 I 讲 5.6，
\[\operatorname{tor.amp}_B\bigl(M_1\mathbin{\overset{\mathrm L}{\otimes}}_{B_1}B\mathbin{\overset{\mathrm L}{\otimes}}_{B_2}M_2\bigr)\subset[a,b],\]
这证明了 b)。另一方面，若 $\operatorname{tor.amp}_{f_1}(E_1)\subset[a_1,b_1]$，则 $M_1$ \Term{isomorphism}{同构}于一个由\Term{flat}{平坦} $B_1$-模组成且在 $[a_1,b_1]$ 外为零的\Term{complex}{复形}。因此，若 $E_2$（从而 $M_2$）在 $[a_2,b_2]$ 外\Term{acyclic}{无上同调}，则 $M_1\mathbin{\overset{\mathrm L}{\otimes}}_{B_1}B\mathbin{\overset{\mathrm L}{\otimes}}_{B_2}M_2$ 在 $[a,b]$ 外\Term{acyclic}{无上同调}，这证明了 a)，也就完成了证明。
\end{proof}
\begin{SGAresult}{推论}{3.5.1}
在 3.5 的假设下，若 $f_1,f_2$ 有\Term{finite}{有限}\Term{tor-dimension}{Tor 维数}（3.1），则 $f$ 也如此。
\end{SGAresult}
\begin{SGAresult}{推论}{3.5.2}
设
\[
\begin{tikzcd}X\arrow[d,"f"']&X'\arrow[l,"g"']\arrow[d,"f'"]\\Y&Y'\arrow[l,"h"]\end{tikzcd}
\]
为\Term{scheme}{概形}的\Term{cartesian-square}{笛卡尔方块}，假定 $X,Y'$ 在 $Y$ 上\Term{tor-independent}{Tor 无关}，并设 $E\in\operatorname{Ob}D(X)$。则
\[\operatorname{tor.amp}_f(E)\subset[a,b]\Longrightarrow\operatorname{tor.amp}_{f'}(\mathrm Lg^*E)\subset[a,b],\]
因此\Term{functor}{函子} $\mathrm Lg^*:D^-(X)\to D^-(X')$ 诱导\Term{functor}{函子}
\[\mathrm Lg^*:D(f)_{\mathrm{torf}}\longrightarrow D(f')_{\mathrm{torf}}.\]
\end{SGAresult}
\SourcePage{256}
\begin{SGAresult}{命题}{3.6}
设
\[
\begin{tikzcd}X\arrow[r,"i"]\arrow[d,"f"']&X'\arrow[dl,"f'"]\\Y&\end{tikzcd}
\]
为\Term{scheme}{概形}的\Term{commutative-diagram}{交换三角形}，其中 $i$ 是\Term{closed-immersion}{闭浸入}，$f'$ 是\Term{smooth}{光滑态射}；设 $E\in\operatorname{Ob}D(X)$。则 $E$ \Term{relative-flat-amplitude}{相对于 $f$ 局部有有限平坦振幅}，当且仅当 $i_*E$（在绝对意义下）\Term{local-scheme}{局部}有\Term{finite}{有限}\Term{flat-amplitude}{平坦振幅}。更精确地，若 $f'$ 的\Term{relative-dimension}{相对维数} $\le d$，则有
\[
\begin{aligned}
\text{(i)}\quad&\operatorname{tor.amp}(i_*E)\subset[a,b]\Longrightarrow\operatorname{tor.amp}_f(E)\subset[a,b],\\
\text{(ii)}\quad&\operatorname{tor.amp}_f(E)\subset[a,b]\Longrightarrow\operatorname{tor.amp}(i_*E)\subset[a-d,b].
\end{aligned}
\]
\end{SGAresult}
\SourcePage{257}
\begin{proof}
若 $[m,n]$ 为 $\mathbf Z$ 中的区间，显然有
\[\operatorname{tor.amp}_f(E)\subset[m,n]\Longleftrightarrow\operatorname{tor.amp}_{f'}(i_*E)\subset[m,n],\]
所以可假定 $X=X'$，$i=\operatorname{Id}_X$，$f=f'$。关系 (i) 直接来自 3.4.1 a)。证明 (ii)。有\Term{cartesian-square}{笛卡尔方块}
\[
\begin{tikzcd}X\arrow[d,"f"']&X\times_YX\arrow[l,"\operatorname{pr}_2"']\arrow[d,"\operatorname{pr}_1"]\\Y&X\arrow[l,"f"]\end{tikzcd}
\]
且\Term{diagonal-morphism}{对角态射} $s:X\to X\times_YX$ 为 $\operatorname{pr}_1,\operatorname{pr}_2$ 的\Term{section}{截面}。假定 $f$ \Term{smooth}{光滑}，\Term{relative-dimension}{相对维数}为 $d$。由第 VII 讲 1.10，$s$ 为\Term{codimension}{余维数} $\le d$ 的\Term{regular-immersion}{正则浸入}，故其\Term{tor-dimension}{Tor 维数} $\le d$（利用 Koszul \Term{complex}{复形}即可看出）。由 3.5.2，
\[\operatorname{tor.amp}_f(E)\subset[a,b]\Longrightarrow\operatorname{tor.amp}_{\operatorname{pr}_1}(\operatorname{pr}_2^*E)\subset[a,b],\]
而由 3.4.1 b)，
\[\operatorname{tor.amp}_{\operatorname{pr}_1}(\operatorname{pr}_2^*E)\subset[a,b]\Longrightarrow\operatorname{tor.amp}(\mathrm Ls^*\operatorname{pr}_2^*E)\subset[a-d,b].\]
由于 $\mathrm Ls^*\operatorname{pr}_2^*\simeq\operatorname{Id}$，蕴含 (ii) 得证。
\end{proof}
\textbf{注记 3.6.1。} 若仅假定 $f'$ \Term{flat}{平坦}，3.6 的陈述便不成立，见 3.4.2。
\begin{SGAresult}{命题}{3.7（\Term{projection-formula}{投影公式}）}
\footnote{参见 [H] II 5.6。}
设 $S$ 为\Term{quasi-compact}{拟紧}\Term{scheme}{概形}，$f:X\to Y$ 为 $S$ 上\Term{quasi-compact}{拟紧}\Term{scheme}{概形}之间的\Term{quasi-compact}{拟紧}、\Term{quasi-separated}{拟分离} $S$-\Term{morphism}{态射}，$E\in\operatorname{Ob}D^b(X)_{\mathrm{qcoh}}$，$M\in\operatorname{Ob}D^b(S)_{\mathrm{qcoh}}$。\footnote{下标“qcoh”表示由具有\Term{quasi-coherent}{拟凝聚}\Term{cohomology}{上同调}的\Term{complex}{复形}组成的\Term{full-subcategory}{满子范畴}。} 若 $M$ 有\Term{finite}{有限}\Term{flat-amplitude}{平坦振幅}，或 $E$ \Term{relative-flat-amplitude}{相对于 $S$ 有限平坦振幅}，则
\SourcePage{258}
存在 $D(Y)_{\mathrm{qcoh}}$ 中\Term{canonical}{典范}的\Term{functorial}{函子性}\Term{isomorphism}{同构}
\[\mathrm Rf_*E\mathbin{\overset{\mathrm L}{\otimes}}_SM\xrightarrow{\sim}\mathrm Rf_*(E\mathbin{\overset{\mathrm L}{\otimes}}_SM)\qquad(\text{见 1.6}).\]
\end{SGAresult}
\begin{proof}
首先注意，存在整数 $N$，使对所有 $i>N$ 和所有\Term{quasi-coherent}{拟凝聚} $\mathcal O_X$-模 $F$，有 $R^if_*F=0$（EGA III 1.4.12、IV 1.7.21）。所以\Term{functor}{函子} $\mathrm Rf_*$ 将 $D^b(X)_{\mathrm{qcoh}}$ 映入 $D^b(Y)$，甚至映入 $D^b(Y)_{\mathrm{qcoh}}$（EGA III 1.4.10、IV 1.7.21）。假定 $E$ \Term{relative-flat-amplitude}{相对于 $f$ 有限平坦振幅}（相应地，$M$ 有\Term{finite}{有限}\Term{flat-amplitude}{平坦振幅}）。\Term{adjunction}{伴随}\Term{arrow}{箭头} $\mathrm Lf^*\mathrm Rf_*E\to E$ 经与 $M$ 作\Term{tensor-product}{张量积}给出 $D^-(X)$ 中的\Term{canonical}{典范}\Term{arrow}{箭头}
\[\mathrm Lf^*(\mathrm Rf_*E\mathbin{\overset{\mathrm L}{\otimes}}_SM)\longrightarrow E\mathbin{\overset{\mathrm L}{\otimes}}_SM,\]
其中 $E\mathbin{\overset{\mathrm L}{\otimes}}_SM\in\operatorname{Ob}D^b(X)$。由\Term{adjunction}{伴随}得到 $D^-(Y)$ 中的\Term{canonical}{典范}\Term{arrow}{箭头}
\[\mathrm Rf_*E\mathbin{\overset{\mathrm L}{\otimes}}_SM\longrightarrow\mathrm Rf_*(E\mathbin{\overset{\mathrm L}{\otimes}}_SM).\tag{*}\]
其源和目标都具有\Term{quasi-coherent}{拟凝聚}\Term{cohomology}{上同调}，因为 $\mathrm Rf_*$ 将 $D^b(X)_{\mathrm{qcoh}}$ 映入 $D^b(Y)_{\mathrm{qcoh}}$（EGA III 1.4.10、IV 1.7.21），而 $\mathbin{\overset{\mathrm L}{\otimes}}_S$ 将 $D^-(X)_{\mathrm{qcoh}}\times D^-(S)_{\mathrm{qcoh}}$ 映入 $D^-(X)_{\mathrm{qcoh}}$（把 $X$ 换成 $Y$ 也一样）。证明 (*) 为\Term{isomorphism}{同构}。问题在 $Y$ 上是\Term{local-scheme}{局部}的，可假定 $S,Y$ 均\Term{affine}{仿射}。在非\Term{lie-bracket}{括号}的情形，逐步归约，立即可化为 $M$ 是普通\Term{quasi-coherent}{拟凝聚} $\mathcal O_S$-模的情形。注意\Term{functor}{函子} $\mathrm Rf_*E\mathbin{\overset{\mathrm L}{\otimes}}_S-$ 和 $\mathrm Rf_*(E\mathbin{\overset{\mathrm L}{\otimes}}_S-)$ 在\Term{category}{范畴} $D^b(S)_{\mathrm{qcoh}}$ 上都是“向\Term{bounded-below}{左有界}的”（way-out à gauche，[H] I 7）。取 $M$ 的\Term{free-module}{自由模}左\Term{resolution}{解消}，再加以\Term{truncation}{截断}，便归结为 $M$ 自由的情形。$R^if_*$ 与\Term{colimit}{归纳极限}可交换，因为 $f$ \Term{quasi-compact}{拟紧}且\Term{quasi-separated}{拟分离}，故甚至可假定 $M$ 是\Term{finite-type}{有限型}\Term{free-module}{自由模}，继而假定 $M=\mathcal O_S$。这时 (*) 按定义就是\Term{isomorphism}{同构}。在\Term{lie-bracket}{括号}中的相应情形，逐步归约，可化为 $M=P^{\sim}$，其中 $P$ 为\Term{flat}{平坦}的 $A$-模，
\SourcePage{259}
$S=\operatorname{Spec}A$。令 $A'=\mathrm D_A(P)$（EGA $0_{\mathrm{IV}}$ 18.2.3）、$S'=\operatorname{Spec}A'$，则只须证明 $\mathrm Rf_*E$ 与\Term{flat}{平坦}\Term{base-change}{基变换} $S'\to S$ 可交换，即\Term{canonical}{典范}\Term{arrow}{箭头}
\[q^*\mathrm Rf_*E\longrightarrow\mathrm Rf'_*p^*E\]
为\Term{isomorphism}{同构}；这里令 $f_{(S')}=f'$，以 $p$（相应地 $q$）表示\Term{canonical-projection}{典范投影} $X_{(S')}\to X$（相应地 $Y_{(S')}\to Y$）。这来自 EGA III 1.4.15、IV 1.7.21。命题 3.7 证毕。
\end{proof}
\begin{SGAresult}{推论}{3.7.1}
设 $S$ 为\Term{scheme}{概形}，$f:X\to Y$ 为 $S$ 上\Term{quasi-compact}{拟紧}\Term{scheme}{概形}之间的\Term{quasi-compact}{拟紧}、\Term{quasi-separated}{拟分离} $S$-\Term{morphism}{态射}，$E\in\operatorname{Ob}D^b(X)_{\mathrm{qcoh}}$。若 $E$ \Term{relative-flat-amplitude}{相对于 $S$ 局部有有限平坦振幅}，则 $\mathrm Rf_*E$ 也如此。更精确地，若 $S$ \Term{quasi-compact}{拟紧}，$N$ 是一个整数，使对每个 $i>N$ 和每个\Term{quasi-coherent}{拟凝聚} $\mathcal O_X$-模 $F$，有 $R^if_*F=0$（EGA III 1.4.12、IV 1.7.21），则
\[\operatorname{tor.amp}_S(E)\subset[a,b]\Longrightarrow\operatorname{tor.amp}_S(\mathrm Rf_*E)\subset[a,b+N].\]
\end{SGAresult}
\begin{proof}
显然只须证明最后一条断言。根据 3.3 的 (i)$\Longleftrightarrow$(iv)，并考虑到 $\mathrm Rf_*$“与\Term{restriction}{限制}到 $S$ 的\Term{open-subset}{开集}可交换”，须证明假设 $\operatorname{tor.amp}_S(E)\subset[a,b]$ 蕴含
\[H^i(\mathrm Rf_*E\mathbin{\overset{\mathrm L}{\otimes}}_SM)=0\]
对每个 $i\notin[a,b+N]$ 和每个\Term{quasi-coherent}{拟凝聚} $\mathcal O_S$-模 $M$ 成立。而上述假设意味着，对所有 $i\notin[a,b]$ 及所有\Term{quasi-coherent}{拟凝聚} $\mathcal O_S$-模 $M$，有 $H^i(E\mathbin{\overset{\mathrm L}{\otimes}}_SM)=0$；利用 3.7 即得所证。
\end{proof}
\SourcePage{260}
\begin{SGAresult}{推论}{3.7.2}
设 $f:X\to Y$ 为\Term{quasi-compact}{拟紧}、\Term{quasi-separated}{拟分离}且\Term{locally-finite}{局部有限}\Term{tor-dimension}{Tor 维数}的\Term{morphism}{态射}（3.1），$E\in\operatorname{Ob}D^b(X)_{\mathrm{qcoh}}$。若 $E$（在绝对意义下）\Term{local-scheme}{局部}有\Term{finite}{有限}\Term{flat-amplitude}{平坦振幅}，则 $\mathrm Rf_*E$ 也如此。
\end{SGAresult}
\begin{proof}合用 3.4.1 a) 与 3.7.1 即可。\end{proof}
